% --- Archon physics formalization source begin ---
% archon:physics
% archon:covers PhyXMiniProblems/problem_phyx_mini_0445.lean
% archon:source-report reports/phyx_mini/problem_phyx_mini_0445.source.json

\chapter{Physics problem phyx\_mini\_0445}
\label{ch:PhyXMiniProblems_problem_phyx_mini_0445}

\paragraph{Problem source.}
\#\# Physical scenario

A container with liquid nitrogen at \$100 \textbackslash{}, \textbackslash{}text\{K\}\$ has a cross-sectional area of \$0.5 \textbackslash{}, \textbackslash{}text\{m\}\textasciicircum{}2\$, as shown in figure. Due to heat transfer, some of the liquid evaporates, and in \$1 \textbackslash{}, \textbackslash{}text\{h\}\$ the liquid level drops \$30 \textbackslash{}, \textbackslash{}text\{mm\}\$. The vapor leaving the container passes through a valve and a heater and exits at \$500 \textbackslash{}, \textbackslash{}text\{kPa\}\$, \$260 \textbackslash{}, \textbackslash{}text\{K\}\$.

\#\# Question

Calculate the volume rate of flow of nitrogen gas exiting the heater.

\#\# Answer choices

- A: A: 0.00352 \textbackslash{}

- B: B: 0.04530 \textbackslash{}

- C: C: 0.01500 \textbackslash{}

- D: D: 0.02526 \textbackslash{}

\#\# Figure caption (auxiliary; use the image as primary evidence)

The image shows a diagram of a system involving liquid nitrogen and a heater. 

- On the left, there is a container divided into two sections. The bottom section contains "Liquid N₂" and is filled with horizontal lines, indicating the liquid phase. The top section, labeled "Vapor," is filled with dots, representing the gaseous phase.

- There is an arrow leading from the vapor section to the right, passing through a circular valve or a symbol that could represent a flow control device.

- This arrow continues to a rectangular box labeled "Heater," which contains a series of looped lines suggesting heating elements.

- An arrow exits from the heater towards the right, indicating the direction of the flow or process.

The diagram likely represents a process where liquid nitrogen is vaporized and then heated as part of a controlled system.

\#\# Dataset metadata

Temperature and Heat Transfer; Multi-Formula Reasoning, Physical Model Grounding Reasoning

\paragraph{Recorded answer/context.}
D: D: 0.02526 \textbackslash{}

\paragraph{Figure/image path.}
/root/proposal\_for\_physic/science-mango/phyx\_mini\_run/phyx\_data/test\_image/445.png

\paragraph{Formalization target.}
create a compiling Lean file with sorry bodies at `PhyXMiniProblems/problem\_phyx\_mini\_0445.lean`. The Lean declarations must preserve the physical quantities, dimensions or dimensional roles, figure labels, governing-law hypotheses, and final relation expressed by this problem.
Use Mathlib/Physlib names found through LeanExplore where available. If a physics API is missing, introduce faithful local abstractions rather than scalar placeholder aliases.

\begin{theorem}[Physics formalization target]
\label{thm:physics:phyx_mini_0445:target}
\lean{PhyXMiniProblems.ProblemPhyXMini0445.outlet_nitrogen_volume_flow_rate}
\uses{def:physics:phyx-mini-0445:phyxminiproblems-problemphyxmini0445-volumeflowincubicmetersperminute, def:physics:phyx-mini-0445:phyxminiproblems-problemphyxmini0445-nitrogenevaporationsetup, def:physics:phyx-mini-0445:phyxminiproblems-problemphyxmini0445-matchesnitrogenevaporationscenario, def:physics:phyx-mini-0445:phyxminiproblems-problemphyxmini0445-matchessuppliednitrogenheaterfigure, def:physics:phyx-mini-0445:phyxminiproblems-problemphyxmini0445-matchesproblemreadouts, def:physics:phyx-mini-0445:phyxminiproblems-problemphyxmini0445-satisfiesnitrogenpropertymodel, def:physics:phyx-mini-0445:phyxminiproblems-problemphyxmini0445-matchesreferencenitrogenpropertydata, def:physics:phyx-mini-0445:phyxminiproblems-problemphyxmini0445-hasphysicalnitrogenflowparameters, def:physics:phyx-mini-0445:phyxminiproblems-problemphyxmini0445-satisfiesevaporationandsteadyflowlaws, def:physics:phyx-mini-0445:phyxminiproblems-problemphyxmini0445-answerchoice, def:physics:phyx-mini-0445:phyxminiproblems-problemphyxmini0445-isreportedoutletflowchoice}
The assigned autoformalize agent should translate this physics problem into Lean declarations in the covered file, with theorem and lemma proof bodies written as `by sorry`.
\end{theorem}
\begin{proof}
This is an autoformalization task, not a proof task. Produce statements that can later be proved by the physics prover without weakening the physical model.
\end{proof}

% --- Archon declaration topology begin ---
\section{Declaration topology}
\begin{definition}[volume Dimension]
  \label{def:physics:phyx-mini-0445:phyxminiproblems-problemphyxmini0445-volumedimension}
  \lean{PhyXMiniProblems.ProblemPhyXMini0445.volumeDimension}
  The physical dimension of volume, L³
\end{definition}
\begin{proof}
  This is the declared model definition of volume Dimension.
\end{proof}
\begin{definition}[specific Volume Dimension]
  \label{def:physics:phyx-mini-0445:phyxminiproblems-problemphyxmini0445-specificvolumedimension}
  \lean{PhyXMiniProblems.ProblemPhyXMini0445.specificVolumeDimension}
  \uses{def:physics:phyx-mini-0445:phyxminiproblems-problemphyxmini0445-volumedimension}
  The physical dimension of thermodynamic specific volume, L³ M⁻¹
\end{definition}
\begin{proof}
  This is the declared model definition of specific Volume Dimension.
\end{proof}
\begin{definition}[mass Flow Rate Dimension]
  \label{def:physics:phyx-mini-0445:phyxminiproblems-problemphyxmini0445-massflowratedimension}
  \lean{PhyXMiniProblems.ProblemPhyXMini0445.massFlowRateDimension}
  The physical dimension of mass flow, M T⁻¹
\end{definition}
\begin{proof}
  This is the declared model definition of mass Flow Rate Dimension.
\end{proof}
\begin{definition}[volume Flow Rate Dimension]
  \label{def:physics:phyx-mini-0445:phyxminiproblems-problemphyxmini0445-volumeflowratedimension}
  \lean{PhyXMiniProblems.ProblemPhyXMini0445.volumeFlowRateDimension}
  \uses{def:physics:phyx-mini-0445:phyxminiproblems-problemphyxmini0445-volumedimension}
  The physical dimension of volume flow, L³ T⁻¹
\end{definition}
\begin{proof}
  This is the declared model definition of volume Flow Rate Dimension.
\end{proof}
\begin{definition}[Length Quantity]
  \label{def:physics:phyx-mini-0445:phyxminiproblems-problemphyxmini0445-lengthquantity}
  \lean{PhyXMiniProblems.ProblemPhyXMini0445.LengthQuantity}
  A nonnegative physical length, used for the observed level drop
\end{definition}
\begin{proof}
  This is the declared model definition of Length Quantity.
\end{proof}
\begin{definition}[Time Quantity]
  \label{def:physics:phyx-mini-0445:phyxminiproblems-problemphyxmini0445-timequantity}
  \lean{PhyXMiniProblems.ProblemPhyXMini0445.TimeQuantity}
  A nonnegative physical duration
\end{definition}
\begin{proof}
  This is the declared model definition of Time Quantity.
\end{proof}
\begin{definition}[Specific Volume Quantity]
  \label{def:physics:phyx-mini-0445:phyxminiproblems-problemphyxmini0445-specificvolumequantity}
  \lean{PhyXMiniProblems.ProblemPhyXMini0445.SpecificVolumeQuantity}
  \uses{def:physics:phyx-mini-0445:phyxminiproblems-problemphyxmini0445-specificvolumedimension}
  A nonnegative thermodynamic specific volume
\end{definition}
\begin{proof}
  This is the declared model definition of Specific Volume Quantity.
\end{proof}
\begin{definition}[Mass Flow Rate Quantity]
  \label{def:physics:phyx-mini-0445:phyxminiproblems-problemphyxmini0445-massflowratequantity}
  \lean{PhyXMiniProblems.ProblemPhyXMini0445.MassFlowRateQuantity}
  \uses{def:physics:phyx-mini-0445:phyxminiproblems-problemphyxmini0445-massflowratedimension}
  A nonnegative mass-flow rate
\end{definition}
\begin{proof}
  This is the declared model definition of Mass Flow Rate Quantity.
\end{proof}
\begin{definition}[Volume Flow Rate Quantity]
  \label{def:physics:phyx-mini-0445:phyxminiproblems-problemphyxmini0445-volumeflowratequantity}
  \lean{PhyXMiniProblems.ProblemPhyXMini0445.VolumeFlowRateQuantity}
  \uses{def:physics:phyx-mini-0445:phyxminiproblems-problemphyxmini0445-volumeflowratedimension}
  A nonnegative volume-flow rate
\end{definition}
\begin{proof}
  This is the declared model definition of Volume Flow Rate Quantity.
\end{proof}
\begin{definition}[area Readout]
  \label{def:physics:phyx-mini-0445:phyxminiproblems-problemphyxmini0445-areareadout}
  \lean{PhyXMiniProblems.ProblemPhyXMini0445.areaReadout}
  Read a physical area in the square of the selected length unit
\end{definition}
\begin{proof}
  This is the declared model definition of area Readout.
\end{proof}
\begin{definition}[length Readout]
  \label{def:physics:phyx-mini-0445:phyxminiproblems-problemphyxmini0445-lengthreadout}
  \lean{PhyXMiniProblems.ProblemPhyXMini0445.lengthReadout}
  \uses{def:physics:phyx-mini-0445:phyxminiproblems-problemphyxmini0445-lengthquantity}
  Read a physical length in the selected length unit
\end{definition}
\begin{proof}
  This is the declared model definition of length Readout.
\end{proof}
\begin{definition}[time Readout]
  \label{def:physics:phyx-mini-0445:phyxminiproblems-problemphyxmini0445-timereadout}
  \lean{PhyXMiniProblems.ProblemPhyXMini0445.timeReadout}
  \uses{def:physics:phyx-mini-0445:phyxminiproblems-problemphyxmini0445-timequantity}
  Read a physical duration in the selected time unit
\end{definition}
\begin{proof}
  This is the declared model definition of time Readout.
\end{proof}
\begin{definition}[specific Volume Readout]
  \label{def:physics:phyx-mini-0445:phyxminiproblems-problemphyxmini0445-specificvolumereadout}
  \lean{PhyXMiniProblems.ProblemPhyXMini0445.specificVolumeReadout}
  \uses{def:physics:phyx-mini-0445:phyxminiproblems-problemphyxmini0445-specificvolumequantity}
  Read a specific volume in cubic selected-length-units per mass unit
\end{definition}
\begin{proof}
  This is the declared model definition of specific Volume Readout.
\end{proof}
\begin{definition}[mass Flow Rate Readout]
  \label{def:physics:phyx-mini-0445:phyxminiproblems-problemphyxmini0445-massflowratereadout}
  \lean{PhyXMiniProblems.ProblemPhyXMini0445.massFlowRateReadout}
  \uses{def:physics:phyx-mini-0445:phyxminiproblems-problemphyxmini0445-massflowratequantity}
  Read a mass-flow rate in selected mass-units per selected time-unit
\end{definition}
\begin{proof}
  This is the declared model definition of mass Flow Rate Readout.
\end{proof}
\begin{definition}[volume Flow Rate Readout]
  \label{def:physics:phyx-mini-0445:phyxminiproblems-problemphyxmini0445-volumeflowratereadout}
  \lean{PhyXMiniProblems.ProblemPhyXMini0445.volumeFlowRateReadout}
  \uses{def:physics:phyx-mini-0445:phyxminiproblems-problemphyxmini0445-volumeflowratequantity}
  Read a volume-flow rate in cubic length-units per selected time-unit
\end{definition}
\begin{proof}
  This is the declared model definition of volume Flow Rate Readout.
\end{proof}
\begin{definition}[temperature In Kelvins]
  \label{def:physics:phyx-mini-0445:phyxminiproblems-problemphyxmini0445-temperatureinkelvins}
  \lean{PhyXMiniProblems.ProblemPhyXMini0445.temperatureInKelvins}
  Read an absolute temperature in kelvins from its zero-preserving unit
\end{definition}
\begin{proof}
  This is the declared model definition of temperature In Kelvins.
\end{proof}
\begin{definition}[pressure In Pascals]
  \label{def:physics:phyx-mini-0445:phyxminiproblems-problemphyxmini0445-pressureinpascals}
  \lean{PhyXMiniProblems.ProblemPhyXMini0445.pressureInPascals}
  Read a physical pressure in coherent SI pascals
\end{definition}
\begin{proof}
  This is the declared model definition of pressure In Pascals.
\end{proof}
\begin{definition}[pressure In Kilopascals]
  \label{def:physics:phyx-mini-0445:phyxminiproblems-problemphyxmini0445-pressureinkilopascals}
  \lean{PhyXMiniProblems.ProblemPhyXMini0445.pressureInKilopascals}
  \uses{def:physics:phyx-mini-0445:phyxminiproblems-problemphyxmini0445-pressureinpascals}
  Read a physical pressure in kilopascals
\end{definition}
\begin{proof}
  This is the declared model definition of pressure In Kilopascals.
\end{proof}
\begin{definition}[volume Flow In Cubic Meters Per Minute]
  \label{def:physics:phyx-mini-0445:phyxminiproblems-problemphyxmini0445-volumeflowincubicmetersperminute}
  \lean{PhyXMiniProblems.ProblemPhyXMini0445.volumeFlowInCubicMetersPerMinute}
  \uses{def:physics:phyx-mini-0445:phyxminiproblems-problemphyxmini0445-volumeflowratequantity, def:physics:phyx-mini-0445:phyxminiproblems-problemphyxmini0445-volumeflowratereadout}
  Cubic metres per minute, the reconstructed unit of the answer list
\end{definition}
\begin{proof}
  This is the declared model definition of volume Flow In Cubic Meters Per Minute.
\end{proof}
\begin{definition}[specific Volume In Cubic Meters Per Kilogram]
  \label{def:physics:phyx-mini-0445:phyxminiproblems-problemphyxmini0445-specificvolumeincubicmetersperkilogram}
  \lean{PhyXMiniProblems.ProblemPhyXMini0445.specificVolumeInCubicMetersPerKilogram}
  \uses{def:physics:phyx-mini-0445:phyxminiproblems-problemphyxmini0445-specificvolumequantity, def:physics:phyx-mini-0445:phyxminiproblems-problemphyxmini0445-specificvolumereadout}
  Cubic metres per kilogram, used for the two nitrogen property entries
\end{definition}
\begin{proof}
  This is the declared model definition of specific Volume In Cubic Meters Per Kilogram.
\end{proof}
\begin{lemma}[length In Meters eq millimeters div thousand]
  \label{lem:physics:phyx-mini-0445:phyxminiproblems-problemphyxmini0445-lengthinmeters-eq-millimeters-div-thousand}
  \lean{PhyXMiniProblems.ProblemPhyXMini0445.lengthInMeters_eq_millimeters_div_thousand}
  \uses{def:physics:phyx-mini-0445:phyxminiproblems-problemphyxmini0445-lengthquantity, def:physics:phyx-mini-0445:phyxminiproblems-problemphyxmini0445-lengthreadout}
  These are general named-unit conversions. They contain no data about the particular nitrogen system and no requested flow-rate value
\end{lemma}
\begin{proof}
  The conclusion follows from the governing relations and figure readouts named in the statement of length In Meters eq millimeters div thousand.
\end{proof}
\begin{lemma}[time In Minutes eq sixty mul hours]
  \label{lem:physics:phyx-mini-0445:phyxminiproblems-problemphyxmini0445-timeinminutes-eq-sixty-mul-hours}
  \lean{PhyXMiniProblems.ProblemPhyXMini0445.timeInMinutes_eq_sixty_mul_hours}
  \uses{def:physics:phyx-mini-0445:phyxminiproblems-problemphyxmini0445-timequantity, def:physics:phyx-mini-0445:phyxminiproblems-problemphyxmini0445-timereadout}
  This lemma records the time In Minutes eq sixty mul hours component of the problem-specific physical model
\end{lemma}
\begin{proof}
  The conclusion follows from the governing relations and figure readouts named in the statement of time In Minutes eq sixty mul hours.
\end{proof}
\begin{definition}[Working Substance]
  \label{def:physics:phyx-mini-0445:phyxminiproblems-problemphyxmini0445-workingsubstance}
  \lean{PhyXMiniProblems.ProblemPhyXMini0445.WorkingSubstance}
  The working substance named in the problem
\end{definition}
\begin{proof}
  This is the declared model definition of Working Substance.
\end{proof}
\begin{definition}[Matter Phase]
  \label{def:physics:phyx-mini-0445:phyxminiproblems-problemphyxmini0445-matterphase}
  \lean{PhyXMiniProblems.ProblemPhyXMini0445.MatterPhase}
  Thermodynamic phases relevant to the liquid container and outlet
\end{definition}
\begin{proof}
  This is the declared model definition of Matter Phase.
\end{proof}
\begin{definition}[Level Drop Cause]
  \label{def:physics:phyx-mini-0445:phyxminiproblems-problemphyxmini0445-leveldropcause}
  \lean{PhyXMiniProblems.ProblemPhyXMini0445.LevelDropCause}
  The mechanism stated to lower the liquid level
\end{definition}
\begin{proof}
  This is the declared model definition of Level Drop Cause.
\end{proof}
\begin{definition}[Container Region]
  \label{def:physics:phyx-mini-0445:phyxminiproblems-problemphyxmini0445-containerregion}
  \lean{PhyXMiniProblems.ProblemPhyXMini0445.ContainerRegion}
  The two labeled regions inside the left-hand container
\end{definition}
\begin{proof}
  This is the declared model definition of Container Region.
\end{proof}
\begin{definition}[Figure Object]
  \label{def:physics:phyx-mini-0445:phyxminiproblems-problemphyxmini0445-figureobject}
  \lean{PhyXMiniProblems.ProblemPhyXMini0445.FigureObject}
  Major objects visible along the left-to-right flow path
\end{definition}
\begin{proof}
  This is the declared model definition of Figure Object.
\end{proof}
\begin{definition}[Figure Text Label]
  \label{def:physics:phyx-mini-0445:phyxminiproblems-problemphyxmini0445-figuretextlabel}
  \lean{PhyXMiniProblems.ProblemPhyXMini0445.FigureTextLabel}
  Literal text labels visible in the supplied raster
\end{definition}
\begin{proof}
  This is the declared model definition of Figure Text Label.
\end{proof}
\begin{definition}[Flow Segment]
  \label{def:physics:phyx-mini-0445:phyxminiproblems-problemphyxmini0445-flowsegment}
  \lean{PhyXMiniProblems.ProblemPhyXMini0445.FlowSegment}
  The three directed flow segments visible in the raster
\end{definition}
\begin{proof}
  This is the declared model definition of Flow Segment.
\end{proof}
\begin{definition}[Flow Node]
  \label{def:physics:phyx-mini-0445:phyxminiproblems-problemphyxmini0445-flownode}
  \lean{PhyXMiniProblems.ProblemPhyXMini0445.FlowNode}
  Nodes joined by the three displayed flow segments
\end{definition}
\begin{proof}
  This is the declared model definition of Flow Node.
\end{proof}
\begin{definition}[expected Segment Start]
  \label{def:physics:phyx-mini-0445:phyxminiproblems-problemphyxmini0445-expectedsegmentstart}
  \lean{PhyXMiniProblems.ProblemPhyXMini0445.expectedSegmentStart}
  \uses{def:physics:phyx-mini-0445:phyxminiproblems-problemphyxmini0445-flowsegment, def:physics:phyx-mini-0445:phyxminiproblems-problemphyxmini0445-flownode}
  Expected initial node of each left-to-right segment
\end{definition}
\begin{proof}
  This is the declared model definition of expected Segment Start.
\end{proof}
\begin{definition}[expected Segment Finish]
  \label{def:physics:phyx-mini-0445:phyxminiproblems-problemphyxmini0445-expectedsegmentfinish}
  \lean{PhyXMiniProblems.ProblemPhyXMini0445.expectedSegmentFinish}
  \uses{def:physics:phyx-mini-0445:phyxminiproblems-problemphyxmini0445-flowsegment, def:physics:phyx-mini-0445:phyxminiproblems-problemphyxmini0445-flownode}
  Expected final node of each left-to-right segment
\end{definition}
\begin{proof}
  This is the declared model definition of expected Segment Finish.
\end{proof}
\begin{definition}[Nitrogen Heater Figure]
  \label{def:physics:phyx-mini-0445:phyxminiproblems-problemphyxmini0445-nitrogenheaterfigure}
  \lean{PhyXMiniProblems.ProblemPhyXMini0445.NitrogenHeaterFigure}
  \uses{def:physics:phyx-mini-0445:phyxminiproblems-problemphyxmini0445-containerregion, def:physics:phyx-mini-0445:phyxminiproblems-problemphyxmini0445-figureobject, def:physics:phyx-mini-0445:phyxminiproblems-problemphyxmini0445-figuretextlabel, def:physics:phyx-mini-0445:phyxminiproblems-problemphyxmini0445-flowsegment, def:physics:phyx-mini-0445:phyxminiproblems-problemphyxmini0445-flownode}
  Qualitative content of the primary image. It records phase-region placement, objects, text, and directed connectivity but contains no numerical flow rate
\end{definition}
\begin{proof}
  This is the declared model definition of Nitrogen Heater Figure.
\end{proof}
\begin{definition}[Liquid Nitrogen State]
  \label{def:physics:phyx-mini-0445:phyxminiproblems-problemphyxmini0445-liquidnitrogenstate}
  \lean{PhyXMiniProblems.ProblemPhyXMini0445.LiquidNitrogenState}
  \uses{def:physics:phyx-mini-0445:phyxminiproblems-problemphyxmini0445-specificvolumequantity, def:physics:phyx-mini-0445:phyxminiproblems-problemphyxmini0445-matterphase}
  The liquid state in the constant-cross-section container
\end{definition}
\begin{proof}
  This is the declared model definition of Liquid Nitrogen State.
\end{proof}
\begin{definition}[Outlet Nitrogen State]
  \label{def:physics:phyx-mini-0445:phyxminiproblems-problemphyxmini0445-outletnitrogenstate}
  \lean{PhyXMiniProblems.ProblemPhyXMini0445.OutletNitrogenState}
  \uses{def:physics:phyx-mini-0445:phyxminiproblems-problemphyxmini0445-specificvolumequantity, def:physics:phyx-mini-0445:phyxminiproblems-problemphyxmini0445-matterphase}
  The nitrogen-gas state immediately after the heater
\end{definition}
\begin{proof}
  This is the declared model definition of Outlet Nitrogen State.
\end{proof}
\begin{definition}[Nitrogen Property Table]
  \label{def:physics:phyx-mini-0445:phyxminiproblems-problemphyxmini0445-nitrogenpropertytable}
  \lean{PhyXMiniProblems.ProblemPhyXMini0445.NitrogenPropertyTable}
  \uses{def:physics:phyx-mini-0445:phyxminiproblems-problemphyxmini0445-specificvolumequantity}
  An external nitrogen property table. The liquid entry is indexed by temperature, while the gas entry is indexed by pressure and temperature
\end{definition}
\begin{proof}
  This is the declared model definition of Nitrogen Property Table.
\end{proof}
\begin{definition}[Nitrogen Evaporation Setup]
  \label{def:physics:phyx-mini-0445:phyxminiproblems-problemphyxmini0445-nitrogenevaporationsetup}
  \lean{PhyXMiniProblems.ProblemPhyXMini0445.NitrogenEvaporationSetup}
  \uses{def:physics:phyx-mini-0445:phyxminiproblems-problemphyxmini0445-lengthquantity, def:physics:phyx-mini-0445:phyxminiproblems-problemphyxmini0445-timequantity, def:physics:phyx-mini-0445:phyxminiproblems-problemphyxmini0445-massflowratequantity, def:physics:phyx-mini-0445:phyxminiproblems-problemphyxmini0445-volumeflowratequantity, def:physics:phyx-mini-0445:phyxminiproblems-problemphyxmini0445-workingsubstance, def:physics:phyx-mini-0445:phyxminiproblems-problemphyxmini0445-leveldropcause, def:physics:phyx-mini-0445:phyxminiproblems-problemphyxmini0445-nitrogenheaterfigure, def:physics:phyx-mini-0445:phyxminiproblems-problemphyxmini0445-liquidnitrogenstate, def:physics:phyx-mini-0445:phyxminiproblems-problemphyxmini0445-outletnitrogenstate, def:physics:phyx-mini-0445:phyxminiproblems-problemphyxmini0445-nitrogenpropertytable}
  Independent physical observables of the experiment. In particular, the outlet volume-flow rate is an independent dimensionful field; it is not defined from an answer choice or from the requested numerical value
\end{definition}
\begin{proof}
  This is the declared model definition of Nitrogen Evaporation Setup.
\end{proof}
\begin{definition}[Matches Nitrogen Evaporation Scenario]
  \label{def:physics:phyx-mini-0445:phyxminiproblems-problemphyxmini0445-matchesnitrogenevaporationscenario}
  \lean{PhyXMiniProblems.ProblemPhyXMini0445.MatchesNitrogenEvaporationScenario}
  \uses{def:physics:phyx-mini-0445:phyxminiproblems-problemphyxmini0445-nitrogenevaporationsetup}
  Qualitative facts stated by the prose scenario
\end{definition}
\begin{proof}
  This is the declared model definition of Matches Nitrogen Evaporation Scenario.
\end{proof}
\begin{definition}[Matches Supplied Nitrogen Heater Figure]
  \label{def:physics:phyx-mini-0445:phyxminiproblems-problemphyxmini0445-matchessuppliednitrogenheaterfigure}
  \lean{PhyXMiniProblems.ProblemPhyXMini0445.MatchesSuppliedNitrogenHeaterFigure}
  \uses{def:physics:phyx-mini-0445:phyxminiproblems-problemphyxmini0445-expectedsegmentstart, def:physics:phyx-mini-0445:phyxminiproblems-problemphyxmini0445-expectedsegmentfinish, def:physics:phyx-mini-0445:phyxminiproblems-problemphyxmini0445-nitrogenheaterfigure}
  Exact qualitative transcription of the supplied raster
\end{definition}
\begin{proof}
  This is the declared model definition of Matches Supplied Nitrogen Heater Figure.
\end{proof}
\begin{definition}[Matches Problem Readouts]
  \label{def:physics:phyx-mini-0445:phyxminiproblems-problemphyxmini0445-matchesproblemreadouts}
  \lean{PhyXMiniProblems.ProblemPhyXMini0445.MatchesProblemReadouts}
  \uses{def:physics:phyx-mini-0445:phyxminiproblems-problemphyxmini0445-areareadout, def:physics:phyx-mini-0445:phyxminiproblems-problemphyxmini0445-lengthreadout, def:physics:phyx-mini-0445:phyxminiproblems-problemphyxmini0445-timereadout, def:physics:phyx-mini-0445:phyxminiproblems-problemphyxmini0445-temperatureinkelvins, def:physics:phyx-mini-0445:phyxminiproblems-problemphyxmini0445-pressureinkilopascals, def:physics:phyx-mini-0445:phyxminiproblems-problemphyxmini0445-nitrogenevaporationsetup}
  Numerical quantities stated in the prose. There is deliberately no numerical readout for the requested outlet volume-flow rate in this structure
\end{definition}
\begin{proof}
  This is the declared model definition of Matches Problem Readouts.
\end{proof}
\begin{definition}[Satisfies Nitrogen Property Model]
  \label{def:physics:phyx-mini-0445:phyxminiproblems-problemphyxmini0445-satisfiesnitrogenpropertymodel}
  \lean{PhyXMiniProblems.ProblemPhyXMini0445.SatisfiesNitrogenPropertyModel}
  \uses{def:physics:phyx-mini-0445:phyxminiproblems-problemphyxmini0445-nitrogenevaporationsetup}
  Constitutive interpretation of the two states: their independent specific volumes are supplied by the nitrogen property table at the measured states
\end{definition}
\begin{proof}
  This is the declared model definition of Satisfies Nitrogen Property Model.
\end{proof}
\begin{definition}[Matches Reference Nitrogen Property Data]
  \label{def:physics:phyx-mini-0445:phyxminiproblems-problemphyxmini0445-matchesreferencenitrogenpropertydata}
  \lean{PhyXMiniProblems.ProblemPhyXMini0445.MatchesReferenceNitrogenPropertyData}
  \uses{def:physics:phyx-mini-0445:phyxminiproblems-problemphyxmini0445-specificvolumeincubicmetersperkilogram, def:physics:phyx-mini-0445:phyxminiproblems-problemphyxmini0445-nitrogenevaporationsetup}
  Reference nitrogen property readouts needed to make the numerical exercise determinate. They are state-property data, not the requested volume-flow rate
\end{definition}
\begin{proof}
  This is the declared model definition of Matches Reference Nitrogen Property Data.
\end{proof}
\begin{definition}[Has Physical Nitrogen Flow Parameters]
  \label{def:physics:phyx-mini-0445:phyxminiproblems-problemphyxmini0445-hasphysicalnitrogenflowparameters}
  \lean{PhyXMiniProblems.ProblemPhyXMini0445.HasPhysicalNitrogenFlowParameters}
  \uses{def:physics:phyx-mini-0445:phyxminiproblems-problemphyxmini0445-areareadout, def:physics:phyx-mini-0445:phyxminiproblems-problemphyxmini0445-lengthreadout, def:physics:phyx-mini-0445:phyxminiproblems-problemphyxmini0445-timereadout, def:physics:phyx-mini-0445:phyxminiproblems-problemphyxmini0445-temperatureinkelvins, def:physics:phyx-mini-0445:phyxminiproblems-problemphyxmini0445-pressureinpascals, def:physics:phyx-mini-0445:phyxminiproblems-problemphyxmini0445-specificvolumeincubicmetersperkilogram, def:physics:phyx-mini-0445:phyxminiproblems-problemphyxmini0445-nitrogenevaporationsetup}
  Positivity conditions selecting the physically meaningful branch
\end{definition}
\begin{proof}
  This is the declared model definition of Has Physical Nitrogen Flow Parameters.
\end{proof}
\begin{definition}[Satisfies Evaporation And Steady Flow Laws]
  \label{def:physics:phyx-mini-0445:phyxminiproblems-problemphyxmini0445-satisfiesevaporationandsteadyflowlaws}
  \lean{PhyXMiniProblems.ProblemPhyXMini0445.SatisfiesEvaporationAndSteadyFlowLaws}
  \uses{def:physics:phyx-mini-0445:phyxminiproblems-problemphyxmini0445-areareadout, def:physics:phyx-mini-0445:phyxminiproblems-problemphyxmini0445-lengthreadout, def:physics:phyx-mini-0445:phyxminiproblems-problemphyxmini0445-timereadout, def:physics:phyx-mini-0445:phyxminiproblems-problemphyxmini0445-specificvolumereadout, def:physics:phyx-mini-0445:phyxminiproblems-problemphyxmini0445-massflowratereadout, def:physics:phyx-mini-0445:phyxminiproblems-problemphyxmini0445-volumeflowratereadout, def:physics:phyx-mini-0445:phyxminiproblems-problemphyxmini0445-nitrogenevaporationsetup}
  The governing flow relations in arbitrary coherent choices of length, mass, and time units: A Δh = m\_dot Δt v\_liquid converts liquid volume loss into evaporated mass; mass flow is conserved through the valve and heater; and V\_dot\_exit = m\_dot\_exit v\_exit converts outlet mass flow to gas volume flow. These general laws contain no answer-choice value and no specialized numerical outlet flow rate
\end{definition}
\begin{proof}
  This is the declared model definition of Satisfies Evaporation And Steady Flow Laws.
\end{proof}
\begin{definition}[Answer Choice]
  \label{def:physics:phyx-mini-0445:phyxminiproblems-problemphyxmini0445-answerchoice}
  \lean{PhyXMiniProblems.ProblemPhyXMini0445.AnswerChoice}
  Labels printed beside the four candidate flow-rate values
\end{definition}
\begin{proof}
  This is the declared model definition of Answer Choice.
\end{proof}
\begin{definition}[displayed Outlet Flow In Cubic Meters Per Minute]
  \label{def:physics:phyx-mini-0445:phyxminiproblems-problemphyxmini0445-displayedoutletflowincubicmetersperminute}
  \lean{PhyXMiniProblems.ProblemPhyXMini0445.displayedOutletFlowInCubicMetersPerMinute}
  \uses{def:physics:phyx-mini-0445:phyxminiproblems-problemphyxmini0445-answerchoice}
  Candidate outlet volume flows, reconstructed as cubic metres per minute
\end{definition}
\begin{proof}
  This is the declared model definition of displayed Outlet Flow In Cubic Meters Per Minute.
\end{proof}
\begin{definition}[recorded Dataset Answer]
  \label{def:physics:phyx-mini-0445:phyxminiproblems-problemphyxmini0445-recordeddatasetanswer}
  \lean{PhyXMiniProblems.ProblemPhyXMini0445.recordedDatasetAnswer}
  \uses{def:physics:phyx-mini-0445:phyxminiproblems-problemphyxmini0445-answerchoice}
  Dataset answer metadata, deliberately not used as a theorem premise
\end{definition}
\begin{proof}
  This is the declared model definition of recorded Dataset Answer.
\end{proof}
\begin{definition}[Is Reported Outlet Flow Choice]
  \label{def:physics:phyx-mini-0445:phyxminiproblems-problemphyxmini0445-isreportedoutletflowchoice}
  \lean{PhyXMiniProblems.ProblemPhyXMini0445.IsReportedOutletFlowChoice}
  \uses{def:physics:phyx-mini-0445:phyxminiproblems-problemphyxmini0445-volumeflowincubicmetersperminute, def:physics:phyx-mini-0445:phyxminiproblems-problemphyxmini0445-nitrogenevaporationsetup, def:physics:phyx-mini-0445:phyxminiproblems-problemphyxmini0445-answerchoice, def:physics:phyx-mini-0445:phyxminiproblems-problemphyxmini0445-displayedoutletflowincubicmetersperminute}
  Agreement at the five decimal places displayed by the answer list
\end{definition}
\begin{proof}
  This is the declared model definition of Is Reported Outlet Flow Choice.
\end{proof}
% --- Archon declaration topology end ---
% --- Archon physics formalization source end ---
